\chapter{原型系统设计与实现}
\label{chap:system}

\section{系统需求分析}
为验证所提框架在实际科研写作场景中的可用性，本文设计并实现了一套协同写作原型系统，涵盖用户管理、文档编辑、文献核验与格式导出等核心模块。

\section{数据库设计}
系统核心用户信息表的字段设计如表~\ref{tab:user}~所示。

\begin{table}[htbp]
  \centering
  \caption{用户信息表字段设计}
  \label{tab:user}
  \begin{tabular}{llll}
    \toprule
    字段名 & 字段代码 & 字段类型 & 备注 \\
    \midrule
    用户 ID & user\_id & INT & 唯一主键（自增） \\
    用户名 & username & VARCHAR & 登录名（唯一） \\
    用户邮箱 & email & VARCHAR & 电子邮箱（唯一） \\
    创建时间 & created\_at & DATETIME & 账号创建时间 \\
    \bottomrule
  \end{tabular}
\end{table}

\section{本章小结}
本章介绍了原型系统的功能模块与核心数据结构设计。
